% Transcription — fonds Alexandre Grothendieck, shelfmark 58, batch 1 (pages 1-20).
% Established from the facsimile archives/batches/58/batch-01.pdf (PDF page k =
% archivists' page k, no cover sheet), cut from a copy of 58.pdf fetched
% through the project's own relay
% (https://grothendieck-relay.onrender.com/source/58.pdf); tiles in
% archives/tiles/58/p1-p20. Per the skill transcribe-grothendieck. The folder
% is sixty-six pages in four batches (1-20, 21-40, 41-60, 61-66), transcribed
% in parallel; this is the first, written without sight of the others.
%
% Pass: Opus 5.5 (claude-opus-5-5), 2026-09-26 — first pass, unchecked against the pages by a human.
%
% Pages skipped: 3, 5, 7, 9, 11, 13, 15, 17 and 19 — the versos, which are
% leaves of a typescript (mostly scanned upside down: graded modules, lambda-
% structures, rational K'-structures on vector spaces) with nothing in his
% hand; skipped by the settled rule for typed versos. Page 1 is a tan cover
% inscribed top right, underlined, « Applicat. à pi. » (the last sign read as
% a pi, doubtful; perhaps pi_1); it takes no \page marker and its words become
% the section heading, with a \note.
%
% Pages summarised: none. Pages transcribed: 2, 4, 6, 8, 10, 12, 14, 16, 18,
% 20 — ten leaves in blue-black ink, fast drafting. The formulas and the
% numbered statements carry; most of the connective prose does not and is
% left \ill{}, as in folder 29 batches 8-9. Pages 14 and 16 are largely
% cancelled by long diagonal strokes and framed blocks; what reads is given
% as \struck{}.
%
% His pagination: none visible on these leaves (the numbers in the lower
% left corner are the archivists'). His numbered runs: Lemme 2 (p. 2),
% Cor. 1 (p. 2), Lemme 3 (p. 4), references to Lemme 1 and Corollaire 2
% (pp. 4, 6; not on these leaves — Lemme 1 must precede or lie in another
% batch), an unnumbered boxed « Théorème » (p. 8), Corollaire 1 (p. 10),
% Corollaire I (p. 12), an unnumbered Corollaire on pi_1 (p. 18), Lemme 4
% (p. 20). Lemme 1 does not occur in pages 1-20.
%
% What the batch is about. Lefschetz-type connectedness for the punctured
% spectrum of a noetherian local ring (the local form of Grothendieck's
% connectedness theorem, SGA 2 Exp. XIII): A local noetherian, X = Spec A,
% X' = X - {m}, Y' = X' cap V(f) for f in the radical rA. Lemme 2: if the
% components of X' have dimension >= 2, Y' is connected. Lemme 3: if the
% components of X' have dim >= k+1 and X' is connected in dimension >= k,
% there is a finite Phi in X' such that V(f) cap Phi empty implies Y'
% connected in dim >= k-1 (using (S_2) and a reference read « Hartshorne »).
% Théorème (p. 8): if the components of Spec(Â) have dim >= k+2 and X' is
% connected in dim >= k, then for f_1..f_l in rA (0 <= l <= k), Y' = X' cap
% V(f_1) cap ... cap V(f_l) satisfies conditions a_{k-l}, b_{k-l}, c_{k-l}
% (connectedness in dim >= k-l), proved by induction on k (Cor. 1, pp. 10-12)
% by choosing f_2 avoiding Phi. Corollaire I (p. 12): a converse — existence
% of f_1..f_k with pi_0 of the intersection controlled forces the dimension
% and connectedness hypotheses on X' (pp. 14-16, largely cancelled). P. 18:
% Corollaire — X' of dim >= n, Y' = X' cap H_1 cap ... cap H_k, k <= n-1:
% Y' connected and pi_1(Y') -> pi_1(X') (surjective, as read), via étale
% covers X'_1 of X'; « Conséquences » for pi_1 of the punctured spectrum.
% P. 20: Lemme 4 — a locally noetherian prescheme whose local punctured
% spectra (Spec O_{X,x})' are connected in dim >= k is connected in dim
% >= k, proved by localising at a point z of X' cap X'' (margin: EGA IV).
%
% Notation fixed once for the batch: \mathfrak{r}A for the radical of A (he
% writes a small r); X' = X - {m}; \hat{A}, \hat{X}; \Phi for the finite set
% of points (a looped capital phi); \pi_0, \pi_1; \underline{\mathcal{O}}_{X,x}
% on p. 20 as he underlines the O; H_i for hyperplane sections on p. 18.
%
% Readings to check first: the cover (« à pi. »); the palimpsest hypotheses
% of Lemme 2 (p. 2); the margin condition « Phi fini » of Lemme 3 (p. 4) and
% its link to « Alors »; « Hartshorne » (p. 6); the definitions a, b, c of
% p. 8; the whole of Corollaire I and pp. 14-16; « surjectif » and the margin
% bracket of p. 18; the superscript of X on p. 20 and the EGA IV margin note.

\documentclass[11pt,a4paper]{article}
% Le préambule commun à toutes les transcriptions du fonds Grothendieck.
%
% Il tient en une idée : **l'incertitude se compose, elle ne se résout pas.**
% Une transcription qui lisse un mot douteux a détruit la seule chose pour
% laquelle on la faisait — un lecteur doit pouvoir savoir, à chaque endroit,
% ce qui était sur le feuillet et ce qui a été ajouté. Les six macros qui
% suivent sont donc l'appareil critique, et non de la décoration.
%
% Les mêmes commandes sont reconnues par `scripts/render.mjs`, qui produit la
% vue de lecture du site. Une macro ajoutée ici doit l'être aussi là-bas,
% faute de quoi le rendu échouera bruyamment — ce qui est voulu : un rendu
% silencieusement faux est pire qu'un rendu absent.

\ProvidesPackage{grothendieck}[2026/08/08 Transcriptions du fonds Grothendieck]

\usepackage{amsmath,amssymb,amsthm}
\usepackage{mathrsfs}
\usepackage{stmaryrd}
% Les diagrammes commutatifs sont partout dans ce fonds : c'est la langue
% même de la géométrie algébrique de Grothendieck, et une transcription qui
% les rendrait en prose ne transcrirait rien.
\usepackage{tikz-cd}
% Français par défaut — c'est la langue des feuillets — mais l'anglais est
% chargé aussi : la traduction et le résumé sont des documents anglais, et la
% typographie française (espaces avant les deux-points) y serait une faute.
% Ces documents basculent par \selectlanguage{english} en tête de corps.
\usepackage[english,french]{babel}
\usepackage{fontspec}
\usepackage{geometry}
\geometry{a4paper,margin=25mm,marginparwidth=28mm}
\usepackage{marginnote}
\usepackage{xcolor}
% ulem, not soul. `soul`'s \st and \ul are built on a character-by-character
% scan that XeLaTeX's font machinery breaks: the document fails with an
% undefined control sequence rather than degrading. ulem does the same two jobs
% and is Unicode-engine safe. `normalem` stops it hijacking \emph, which the
% transcriptions use for Grothendieck's own emphasis.
\usepackage[normalem]{ulem}
\usepackage{hyperref}
\hypersetup{colorlinks=true,linkcolor=black,urlcolor=black,pdfborder={0 0 0}}

% --- Métadonnées du lot -----------------------------------------------------
% Reprises telles quelles par le script de rendu, qui les affiche en tête de
% la vue de lecture. Elles fixent ce que le fichier prétend couvrir : sans
% elles, une transcription flotte sans point d'attache dans un fonds de seize
% mille feuillets.

\newcommand{\folder}[1]{\gdef\grothFolder{#1}}
\newcommand{\batch}[1]{\gdef\grothBatch{#1}}
\newcommand{\leaves}[2]{\gdef\grothFirst{#1}\gdef\grothLast{#2}}
\newcommand{\foldertitle}[1]{\gdef\grothFoldertitle{#1}}
\gdef\grothFolder{?}\gdef\grothBatch{?}\gdef\grothFirst{?}\gdef\grothLast{?}\gdef\grothFoldertitle{}

% --- Le résumé --------------------------------------------------------------
%
% Il ouvre la lecture modernisée, sous le titre « Résumé », et il est composé
% en petit. Ce n'est pas un ornement : c'est le seul passage du document qui
% ne soit pas des mathématiques, et un lecteur doit pouvoir le reconnaître
% avant de l'avoir lu — puis le sauter s'il connaît déjà le sujet, ou s'y
% arrêter s'il ne le connaît pas. Le filet qui le suit marque la reprise du
% corps.

\newenvironment{resume}
  {\small\setlength{\parskip}{0.45em}}
  {\par\medskip\hrule\medskip}

% --- Mots-clés ---------------------------------------------------------------
%
% En anglais, en clôture du résumé de la lecture modernisée : le vocabulaire
% moderne sous lequel chercher ce que les feuillets construisent. Le manifeste
% du site (`npm run manifest`) les extrait pour étiqueter la cote entière —
% cette ligne est la source unique des tags, il n'y a pas de fichier à côté.
\newcommand{\keywords}[1]{\par\smallskip\noindent
  {\footnotesize\textsc{Keywords} --- \textit{#1}}\par}

% --- L'appareil critique ----------------------------------------------------

% Le numéro de feuillet, en marge. C'est l'ancre qui permet de régler un
% désaccord sur une formule en regardant *un* feuillet plutôt qu'une chemise
% de six cents. Le site s'en sert aussi pour tourner les pages du fac-similé
% quand on fait défiler la transcription.
\newcommand{\leaf}[1]{%
  \marginnote{\footnotesize\color{gray!70}\textsf{#1}}%
  \label{leaf:#1}\ignorespaces}

% Illisible : on ne devine pas. Un mot inventé qui se lit comme les autres est
% le pire résultat possible.
\newcommand{\ill}{\textcolor{red!60!black}{[\,\ldots\,]}}

% Lecture douteuse : le mot est donné, et signalé comme incertain.
\newcommand{\uncertain}[1]{\uline{#1}}

% Ajout de l'éditeur — une lettre manquante, un mot sous-entendu.
\newcommand{\add}[1]{\textcolor{blue!50!black}{[#1]}}

% Biffé par Grothendieck lui-même. Ce qu'il a rayé fait partie du document :
% une rature dit souvent où la pensée a changé de direction.
\newcommand{\struck}[1]{\sout{#1}}

% Note du transcripteur, sur l'état matériel du feuillet.
\newcommand{\note}[1]{\par\smallskip{\small\color{gray!60!black}\textit{[#1]}}\par\smallskip}

% Note marginale de Grothendieck, distincte du corps du texte.
\newcommand{\marginal}[1]{\par\smallskip\noindent\hspace{1em}%
  {\small\color{gray!60!black}#1}\par\smallskip}

% --- Titre ------------------------------------------------------------------

\AtBeginDocument{%
  \begin{center}
    {\large\bfseries Fonds Alexandre Grothendieck --- cote n\textsuperscript{o}~\grothFolder}\\[2pt]
    {\itshape\grothFoldertitle}\\[4pt]
    {\small feuillets \grothFirst--\grothLast\ (lot \grothBatch)}\\[2pt]
    {\footnotesize\color{gray!60!black}Transcription établie d'après le fac-similé numérisé
     par l'Université de Montpellier.\\
     Les lectures incertaines sont soulignées, les illisibles notées [\,\ldots\,],
     les ajouts entre crochets.}
  \end{center}
  \vspace{1em}}

\endinput


\folder{58}
\batch{1}
\pages{1}{20}
\dating{[à partir de 1964]}
\watermark{Édition de démonstration}
\foldertitle{Théorie de " Lefschetz-Grauert ". Applications en π₁, à Pic etc… : notes manuscrites (s.d.)}

\begin{document}

\section{Applicat. à $\pi$.}
\note{Titre écrit par Grothendieck, souligné, en haut à droite de la couverture (page 1) ; le dernier signe, lu comme un $\pi$, est d'une lecture incertaine (peut-être $\pi_1$).}

\page{2}
\struck{Les composantes de $X$ de dim $\geq 3$}

\emph{Lemme 2.} $A$ anneau local noethérien, $f \in \mathfrak{r}A$, $X = \operatorname{Spec} A$, $X' = X - \lbrace \mathfrak{m} \rbrace$, $Y' = X' \cap V(f)$.
Supposons \ill{} \struck{les} composantes \ill{} de $X'$, \struck{\ill{}} \struck{\ill{} de $X'$ de dim $\geq 2$} \struck{\ill{}} \struck{\ill{} de dim $\geq 1$} \struck{\ill{}} d'un \ill{} \ill{} de $X'$ \ill{} de dim $\geq 2$. Alors \ill{} $Y'$ \ill{} de dim $\geq 1$, \uncertain{et} $Y'$ est \emph{\uncertain{connexe}}.
\marginal{complet}
\note{le mot « complet », encerclé en haut à droite, est relié par un trait à « noethérien » : c'est un ajout à l'énoncé.}
\note{$\mathfrak{r}A$ : le radical de $A$ ; la lettre est lue d'après le contexte.}
\note{Le milieu de l'énoncé est un palimpseste de quatre lignes biffées et récrites dans l'interligne ; seule la forme finale des hypothèses (composantes de $X'$ de dimension $\geq 2$) se lit avec quelque assurance.}

En effet, \ill{} \ill{} $Z'$  de $Y'$, \struck{l'intersection \ill{} de dim \ill{}} les composantes \ill{} $Z' \cap Y'$ \ill{} de dim $\geq \dim Z' - 1$ ;  \ill{} \ill{} (si $Z' = X'$) \ill{} $Z'$ est de dim $\geq 1$. De plus, \ill{} pour toute partie \ill{} de dim $\geq 0$. Donc \ill{} si $Z'$ de dim $\geq 1$, $Z' \cap Y'$ \ill{} \ill{} \ill{}.
\marginal{fermée-ouverte}
\marginal{si $Z' \cap Y' \neq \emptyset$ si $\dim Z' > 0$}
\note{« fermée-ouverte » : mot encerclé en marge, relié à la ligne, lecture incertaine ; la seconde note marginale est reliée à « de dim $\geq \dim Z' - 1$ ».}

\emph{Cor. 1.} \ill{} Lemme 2, \ill{} \ill{} \ill{} \ill{}, $X'$ \ill{} irréductible, et \ill{} \ill{} \ill{} \ill{} $X$ intègre \ill{} de dim $\geq 2$. [D'ailleurs \ill{}, \ill{} \ill{} \ill{} \ill{} \ill{} \ill{} \ill{} $A$ \ill{} \emph{complet}, \ill{} l'hypothèse \ill{} \ill{} \ill{} complet. N.B. \ill{} \ill{} \ill{}, \ill{} \ill{} \ill{} \ill{} \ill{}] Alors $\tilde{A}$ clôture normale \ill{}

\page{4}
de $A$ \ill{} \uncertain{fini} \ill{} $A$, et \ill{} \uncertain{intègre} \ill{} \ill{} \ill{} local. Appliquer \ill{} \uncertain{Lem.~2} \ill{}, \ill{} \ill{} \ill{} \ill{} \ill{} $A$ local \emph{\ill{}} (\ill{} \uncertain{normal}) \ill{} de dim $\geq 3$, \ill{} $X'$ \ill{} \ill{} de dim $\geq 2$. Alors
\[
  H^0(\hat{X}', \mathcal{O}_{\hat{X}'}) \simeq H^0(X', \mathcal{O}_{X'}) \quad (\ill{}) \ \ill{}
\]
\uncertain{conclusion}.
\note{Les indices et l'accent de la formule sont lus d'après le tracé ; le signe entre les deux membres est une flèche double doublée d'un tilde, rendue par $\simeq$.}

\emph{Lemme 3.} $A$ anneau local noethérien \struck{\ill{}} \ill{}, \struck{$f \in \mathfrak{r}A$, \ill{} \ill{}} \struck{\ill{}} $X$, $X'$, \struck{$Y' = V(f) \cap X'$}. \ill{} $k \geq 1$ tel que
\begin{itemize}
  \item[(i)] les composantes \ill{} de $X'$ \ill{} de dim $\geq k+1$ ;
  \item[(ii)] $X'$ est connexe en dim $\geq k$.
\end{itemize}

Alors \struck{$\exists\, f \in \mathfrak{r}A$ \ill{} $\Rightarrow$} $Y' = X' \cap V(f)$ \struck{\ill{} connexe en dim $\geq k$,} \ill{} \ill{} dim $\geq k-1$.
\marginal{$k \geq 1$ tel}
\marginal{$\exists\, \Phi$ fini en $X'$ \ill{} \ill{} $f \in \mathfrak{r}A$, $V(f) \cap \Phi = \emptyset$}
\note{La marge relie par un trait sa condition « il existe $\Phi$ fini en $X'$ … » au « Alors » ; l'énoncé semble donc : pour tout $f \in \mathfrak{r}A$ tel que $V(f)\cap\Phi = \emptyset$, $Y'$ est connexe en dimension $\geq k-1$.}

\uncertain{De façon} \ill{}, \ill{} \ill{} \ill{} \ill{} \ill{} $\Phi$ de $X'$ \ill{} \ill{} \ill{} $V(f) \cap \Phi = \emptyset \Longrightarrow$ \uncertain{conclusion} valable.
Il suffit \ill{} \ill{} \ill{} 1 : \uncertain{Lemme}~1 \ill{} $V(f) \cap X'$ \ill{} connexe en dim $\geq k-1$.

\page{6}
Cela nous ramène au cas où $X$ \uncertain{irréductible}. \ill{} \ill{} \uncertain{Corollaire}~2 \ill{} \ill{} Lemme~1, \ill{} \ill{} \ill{} \ill{} \ill{} \ill{} \ill{} $X$ \ill{} de plus \emph{\uncertain{normal}}, et \ill{}. \uncertain{Soit} $\Phi$ \ill{} \ill{} \ill{}, \ill{} \ill{} \emph{\ill{} \ill{} \ill{}} \emph{\ill{}} $\Phi$ \ill{} \ill{} de $X'$ telle que $V(f) \cap \Phi = \emptyset$ \ill{} $V(f) \cap X' = Y'$ \ill{} \ill{} \ill{} \ill{} $(S_2)$, [\ill{} \ill{} \ill{} $(S_2)$, \ill{} \ill{} \ill{} \ill{} \ill{} $Z'_i$ \ill{} $X'$ ; \ill{} \ill{} \ill{} \ill{} $\Phi$ \ill{} \ill{} \ill{} \ill{} \ill{} \struck{\ill{}} $Z'_i$]. Mais alors, \uncertain{(Hartshorne)} $Y'$ \ill{} \ill{} \ill{} \ill{} \ill{} \ill{} 1, \ill{} \ill{} \ill{} \ill{} \ill{} \ill{} \ill{} \ill{} $\geq k-1$.
\note{Page presque entièrement en écriture rapide ; seules les formules, la condition $(S_2)$ et le renvoi à Hartshorne (lecture incertaine) se lisent. Le raisonnement semble réduire au cas $X$ normal irréductible, choisir un ensemble fini $\Phi$ de points de $X'$ évité par $V(f)$, et conclure à la connexité de $Y'$ en dimension $\geq k-1$.}

\page{8}
\emph{Théorème.} $A$ anneau local noethérien ($X$, $X'$, $k \geq 1$). Supposons \struck{\ill{} \ill{}} les composantes \uncertain{irréductibles} de $\hat{X} = \operatorname{Spec}(\hat{A})$ \ill{} de dim $\geq k+2$ [il suffit pour ceci que celles de $X = \operatorname{Spec} A$ \ill{} de dim $\geq k+2$, i.e.\ celles de $X'$ de dim $\geq k+1$, et \ill{} \ill{} $A \to$ \ill{} \ill{}]. \struck{Soient $f_i$} \ill{} $X'$ est connexe en dim $\geq k$.
\marginal{Conditions \ill{} \ill{} \ill{} $\hat{A}$, \ill{} \ill{} \ill{} \ill{} \ill{} $A$ [\ill{} \ill{} \ill{} \ill{} ici]}
\note{Le mot « Théorème » est encadré.}
\note{La note « Conditions … ici », oblique en marge gauche, est reliée par un crochet aux hypothèses du théorème ; presque illisible.}

Considérons \struck{\ill{} \ill{}} $0 \leq l \leq k$, et $f_1, \ldots, f_l \in \mathfrak{r}A$, \struck{$\bigcap V(f_i)$} $Y' = X' \cap \bigcap_i V(f_i)$. Alors les conditions \ill{} $a_{k-l}$, $b_{k-l}$, \ill{} \ill{} \struck{\ill{}} $Y'$ \ill{} \ill{} \ill{} en dim $\geq k-l$ \ill{} $c_{k-l}$ \ill{} \ill{} \ill{}. \uncertain{Remarque} \ill{} \ill{} \ill{} \ill{} \ill{}.
\note{Les conditions $a_{k-l}$, $b_{k-l}$, $c_{k-l}$ ne sont pas définies sur cette page.}

\emph{Dém.} Si \ill{} \ill{} \ill{} \ill{} $Y'$ \ill{} $Z'$ \ill{} dim $< k-l$, \uncertain{il existe} \ill{} \ill{} \ill{} $\mathfrak{r}A$, \ill{} \ill{} existence ($f_{l+1}, \ldots, f_k$ tels que
\[
  Z' \cap V(f_{l+1}) \cap \cdots \cap V(f_k) = \emptyset,
\]
\ill{} \ill{} \ill{}

\page{10}
\ill{} $Y' \cap V(f_{l+1}) \cap \cdots \cap V(f_k) = X' \cap V(f_1) \cap \cdots \cap V(f_k)$ \ill{} \uncertain{discontinu}\ill{}, \ill{} \ill{} \ill{}.

\emph{Corollaire 1.} Sous les conditions \ill{} \ill{} \ill{} \ill{} \ill{} ($X$, $X'$, $k$), \struck{Soient} \ill{} $f_1, \ldots, f_k \in \mathfrak{r}A$, \ill{}
\[
  Y' = V(f_1) \cap \cdots \cap V(f_k) \cap X'
\]
\ill{} \ill{}. \uncertain{On peut} \ill{} \ill{} \ill{} \ill{} \ill{} \ill{} $A$ \ill{} \emph{\ill{}}, \ill{} \ill{} \ill{} \ill{}, \ill{}\ill{}. \uncertain{Dém.} \ill{} \ill{}.
\note{Tout ce passage est rapide ; « Corollaire 1 » est souligné, l'énoncé n'est lisible que dans ses formules.}

Pour $k = 1$, c'est le Lemme~2. \uncertain{Supposons} $k \geq 2$, et le th.\ \uncertain{prouvé} \ill{} \uncertain{dimension} pour les $k' < k$. \uncertain{Le} \ill{} \ill{} : \ill{} \ill{} \ill{} \ill{} \ill{} $k$, et $l = 1$. \uncertain{Procédons} \ill{}\ill{} : \ill{}, \ill{} \ill{} \ill{} \ill{} $Y' = V(f_1) \cap X'$ \ill{} \ill{} \ill{} \ill{} $Z'$ \ill{} $Z'$ de dim $< k-1$, donc \ill{} $\exists\, f_2, \ldots, f_k \in \mathfrak{r}A$ tels que $V(f_1) \cap Z' \cap \cdots \cap V(f_k) = \emptyset$, et \ill{}

\page{12}
\ill{} les \uncertain{systèmes} $X' \cap V(f_1) \cap \cdots \cap V(f_k)$ \ill{} \emph{\uncertain{discontinus}}. Or \struck{\ill{}} \ill{} \ill{} \ill{} Lemme~3, $\exists\, \Phi \subset X'$ fini \ill{} \ill{} \ill{} $f \in \mathfrak{r}A$, $V(f) \cap \Phi = \emptyset \Longrightarrow$ \ill{} composantes \ill{} dim $\leq k-1$. \struck{\ill{} \ill{} $\Phi$ \ill{} \ill{}} \uncertain{Choisissons} les $f_i$ tels que $V(f_2) \cap \Phi = \emptyset$, \ill{} $f_2$ \ill{}, \ill{} \ill{} [\ill{} \ill{} \ill{} \ill{} \ill{} \ill{} $\Phi$ \ill{} \ill{} \ill{} \ill{} les \ill{} \struck{\ill{}} \ill{} \ill{} $Z'$]. Alors $X' \cap V(f_2)$ \ill{} \ill{} \ill{} ($a_{k-1}$, $b_{k-1}$). En \uncertain{vertu} \ill{} \ill{} \ill{} \ill{} \ill{} \ill{}, \ill{} $(X' \cap V(f_2)) \cap V(f_1) \cap V(f_3) \cap \cdots \cap V(f_k)$ \ill{} \ill{} \ill{} \ill{} \emph{\ill{}}.
\note{« Or » est souligné trois fois.}



\emph{Corollaire I.} $(k \geq 1)$ \ill{} \ill{} \ill{}, \ill{} \ill{} \ill{} \ill{} $\mathfrak{r}A$, \ill{} \ill{} \ill{} \ill{} $f_1, \ldots, f_k$ d'\uncertain{idéal} \ill{} \ill{} \ill{} \ill{} $\pi_0(X' \cap V(f_1) \cap \cdots \cap V(f_k))$ \ill{}. \struck{\ill{} \ill{} \ill{} $k-1$}, Alors \uncertain{si} les composantes de $X'$ \ill{} de dim $\geq k+1$, et $X'$ \ill{} connexe en dim $\geq k$.
\marginal{\uncertain{Exclure} \ill{} \ill{} $X$ \uncertain{irréductible} \ill{} \struck{Exclure \ill{} \ill{}} $k \geq 1$ [\struck{$X$ irréductible} \ill{} de dim $1$,}
\marginal{\uncertain{Remarque}}
\note{La note « Exclure … de dim 1 » est en marge gauche, en bas de page, en partie biffée ; lecture très incertaine.}
\note{« Remarque » : mot en marge, souligné, relié par une flèche au titre « Corollaire I ».}
\note{« I » souligné à part, et « $(k \geq 1)$ » ajouté au-dessus par un trait.}

\page{14}

En effet, \ill{} \ill{} $X' \neq \emptyset$. Soit $X'_0$ \ill{} \ill{} \ill{} \ill{} \ill{} \ill{} \ill{}, \ill{} $l$, \ill{} \ill{} ($l \geq 0$), \ill{} $l \geq k+1$. \struck{\ill{} \ill{}}
\note{Page en grande partie annulée par de longs traits obliques et des cadres barrés ; on donne ce qui se lit, les passages annulés étant marqués comme biffés.}

\struck{Existence \ill{} $f_1, \ldots, f_l$ si \ill{} \ill{} \ill{} $l \leq k$, \ill{} $X'$ \ill{} \ill{} \ill{} \ill{} (\ill{} \ill{} $l = 0$] \ill{} \ill{} \ill{} \ill{} \ill{} $X' = X'_0$, \ill{} $X'$ de dim $0$, \ill{} \ill{} \ill{} $\exists\, f \in \mathfrak{r}A$ \ill{} $Y' = X' \cap V(f) = \emptyset$, \ill{}. \ill{} \ill{} $X' \cap V(f_1) \cap \cdots \cap V(f_l) = \emptyset$ \ill{} \ill{}. \ill{} \ill{} $l \leq k$, \ill{} \ill{}.}

\ill{} \ill{} \ill{} $\mathfrak{r}A$ \struck{\ill{}} $f_1, \ldots, f_l$ \ill{} \ill{} \ill{} \ill{} $X'_i$ de dim $n_i$, \ill{} $V(f_1) \cap \cdots \cap V(f_l) \cap X'_i = \emptyset$ \struck{\ill{}} \ill{} $n_i - l$, et \struck{\ill{}} \ill{} \ill{} \uncertain{identique} \ill{} \struck{\ill{}} ($i = 1, \ldots, n$).

\struck{\ill{} \ill{} \ill{} $X'_i \cap V(f_1) \cap \cdots \cap V(f_l)$. \ill{} \ill{} \ill{} \ill{} \ill{}, \ill{} \ill{} \ill{} de dim $0$ et \ill{} \ill{} $l \leq k$,} \ill{} \ill{} \ill{}, \ill{} $l = 0$, $Y'$ \ill{} \ill{} ; \ill{} \ill{} \ill{} $Y'$ \ill{} \ill{} ($l \leq k$, \struck{\ill{}}) \ill{} \ill{} \ill{} \ill{} $Y'$ \ill{} \ill{}, \ill{} \ill{} \ill{} \ill{} \ill{} \ill{} \ill{} \ill{} \ill{} \ill{} $l \leq k$ ; $X'$ \ill{} \ill{} \struck{$l \leq k$, \ill{}}

\page{16}

\struck{Plus \uncertain{précisément} \ill{} $X$ vide si $l = k$, \ill{} \ill{} $X' = k$ \ill{} \ill{} \ill{} $X = k+1 \geq l$, \ill{} \ill{} \ill{} $f_l, \ldots, f_k$, $X'' = X' \cap V(f_l) \cap \cdots \cap V(f_k)$ \ill{} \ill{} $l = 1$ \ill{} $X = 2$ \ill{} $X$ \ill{} de dim $l+1$ \ill{} \ill{} $l = 0$, i.e.\ \ill{} \ill{} de dim $1$.}
\note{Le haut de la page est enfermé dans un cadre barré de traits obliques ; on le donne biffé, dans la mesure où il se lit.}

Nous avons \ill{} \ill{} \ill{} $X$ \ill{} \ill{} \ill{} \ill{} de dim $\geq 2$, \ill{} \ill{} \ill{} \ill{} $f_1, \ldots, f_l$ de $\mathfrak{r}A$ tels que $X' \cap V(f_1) \cap \cdots \cap V(f_l)$ \uncertain{discontinu}. On \ill{} il existe $f_1, \ldots, f_{l-1}$ tels que $Y' = X' \cap V(f_1) \cap \cdots \cap V(f_{l-1})$ \ill{} \ill{} de dim $1$, \ill{} \ill{} \ill{} \ill{} $Y'$ \ill{} \ill{} de dim $2$. Mais alors \ill{} il est \ill{} \ill{} \ill{} $f_l \in \mathfrak{r}B$ tels que $V(f_l)$ \ill{} $Y'$, [\ill{} \ill{} \ill{} de dim $0$] \ill{} \ill{} \ill{} \ill{} \ill{} : \ill{} \ill{} \ill{} \ill{} $Y'$ \ill{} \ill{} \ill{} \ill{} \ill{} \ill{} \ill{}\,!!
\note{Dans « de dim $\geq 2$ », « $l+1$ » est écrit au-dessus du « 2 », en correction.}

\page{18}
\emph{Corollaire.} $X'$ \ill{} \ill{} de dim $\geq n$,
$Y' = X' \cap H_1 \cap \cdots \cap H_k$, \ill{} $k \leq n-1$. Alors \struck{\ill{}} $Y' \to$ \ill{} \ill{} \ill{} et $\pi_1(Y') \to \pi_1(X')$ \ill{} \ill{} : \ill{} \ill{}.
\marginal{[\ill{} \ill{} \ill{}] \uncertain{connexe}}
\note{Ajout encadré par un trait vertical, en haut à droite ; « connexe », écrit plus petit en bout de ligne, est d'une lecture incertaine.}

En effet, le \ill{} \ill{} \ill{} \ill{} $\pi_1$, il \ill{} \ill{} \ill{} \ill{} \ill{} \ill{} $X'_1$ \ill{} \ill{} \ill{} \ill{}, \ill{} \ill{} \ill{} $X'$ \ill{} \emph{\ill{}} \ill{} $X'$, alors $Y'_1 \to$ \ill{}. Or $X'_1$ \ill{} \emph{\uncertain{irréductible}}, \ill{} \ill{} \ill{} \ill{} \ill{}, \ill{} \ill{} \ill{}.

\emph{\struck{\ill{}} Conséquences} \ill{} \ill{} : \ill{} l'\uncertain{ingrédient} \ill{} \ill{} \ill{} $\pi_1(X')$ \ill{} \ill{} \ill{} local $A$, \ill{} \ill{} \ill{} \ill{} \ill{} \ill{} \ill{} \ill{} [\ill{} \ill{}] : \ill{} \ill{} \ill{} \ill{} \ill{} \ill{} \ill{} \ill{} \ill{} \ill{} \ill{} de dim $2$ \ill{} \ill{}. \struck{\ill{} \ill{}} \ill{} \ill{} \ill{} \ill{} \ill{} \ill{} \ill{} \ill{} \ill{} \ill{} \emph{\ill{}} \ill{} \ill{} \ill{} \ill{} \ill{}\,!\ \ill{}.
\note{« Conséquences » remplace un mot biffé et souligné.}

\page{20}
\emph{Lemme 4.} Soit $X$ un \uncertain{préschéma} \ill{}, dont les composantes irréductibles \ill{} de dim $\geq k$, \ill{} \ill{} \uncertain{localement} \ill{} [\ill{} \ill{} \ill{} \ill{} $\mathfrak{r}A$\ill{}], \ill{} \uncertain{tel que} pour $x \in X$, $(\operatorname{Spec} \underline{\mathcal{O}}_{X,x})'$ soit connexe en dim $\geq k$\uncertain{.} Alors $X$ est connexe en dim $\geq k$.
\marginal{\uncertain{Pour} \ill{} \ill{} \ill{} (\ill{} $\mathcal{O}_{X,x}$ \ill{}) \ill{} \ill{}. Alors il \ill{} \ill{} \ill{} \ill{} \ill{} \ill{} \ill{} \ill{} \ill{} \ill{} \ill{} \ill{} \ill{} \ill{} \ill{} EGA IV \ill{}.}
\note{Ajout interligne, souligné : « dont les composantes irréductibles … de dim $\geq k$ ».}
\note{« Pour … EGA IV » : longue note oblique en marge gauche, soulignée par endroits ; on n'en lit guère que le renvoi à EGA~IV.}

\struck{\uncertain{Supposons}} \ill{} \ill{} \ill{}, \ill{} \ill{} \ill{} \ill{} \ill{}. Soit \ill{} \struck{\ill{}} $Z$ \ill{} \ill{} \ill{} $X$ de dim $< k$, \ill{} \ill{} \ill{} \ill{} \ill{} $Z$ ne \uncertain{disconnecte} \ill{} $X$. \struck{\ill{} \ill{}} Le \ill{} \ill{} \ill{} \uncertain{disconnectant}, \ill{} \ill{} \ill{} $U = X - Z = U' \cup U''$, $U', U''$ fermés dans $U$, $U' \cap U'' = \emptyset$. Soient $X' = \overline{U'}$, $X'' = \overline{U''}$, \ill{} \ill{} \uncertain{composantes} $U$ dans \struck{$X = (Z \cup X') \cup X''$, \ill{} donc $(Z \cup X') \cap X'' \neq \emptyset$} $X$, $X^{z} = X' \cup X''$, donc $X' \cap X'' \neq \emptyset$. Soit $z$ \ill{} \ill{} \ill{} de $X' \cap X''$. Alors \uncertain{d'ici} \ill{} \ill{} \ill{} $\operatorname{Spec} \mathcal{O}_{X,z} = X_z$, \ill{} \ill{} \ill{} \ill{} $Z_z$ \uncertain{disconnecte} $X_z$, \ill{} \ill{} \uncertain{disconnecte} $X'_z$, \uncertain{absurde}\ldots
\note{Dans $X^{z} = X' \cup X''$, l'exposant de $X$ est d'une lecture incertaine.}

\end{document}
